\documentclass[11pt]{article}
\usepackage[T1]{fontenc}
\usepackage{lmodern}
\usepackage{amsmath,amssymb,amsthm}
\usepackage[margin=1in]{geometry}
\usepackage{microtype,booktabs,array,longtable}
\usepackage{xcolor}
\usepackage{hyperref}
\usepackage{listings}
\usepackage{fancyhdr}
\pdfmapfile{+lm.map}
\pdfmapfile{+cm.map}
\pdfmapfile{+symbols.map}
\hypersetup{colorlinks=true,linkcolor=blue!45!black,urlcolor=blue!45!black,citecolor=blue!45!black,
 pdftitle={A one-body context calculus for intuitionistic primitive recursive arithmetic},
 pdfauthor={Research prepared for Vladimir Reshetnikov}}
\pagestyle{fancy}
\fancyhf{}
\fancyhead[L]{\small One-body contexts for iPRA}
\fancyhead[R]{\small Research report}
\fancyfoot[C]{\thepage}
\setlength{\headheight}{14pt}
\setlength{\parskip}{4pt}
\setlength{\parindent}{0pt}
\emergencystretch=2em
\newtheorem{theorem}{Theorem}[section]
\newtheorem{lemma}[theorem]{Lemma}
\newtheorem{corollary}[theorem]{Corollary}
\theoremstyle{definition}
\newtheorem{definition}[theorem]{Definition}
\theoremstyle{remark}
\newtheorem{remark}[theorem]{Remark}
\newcommand{\ipra}{\mathrm{iPRA}}
\newcommand{\Bprf}{\mathsf{B}}
\newcommand{\Cprf}{\mathsf{C}}
\newcommand{\Nprf}{\mathsf{N}}
\newcommand{\Teq}{T_{\mathrm{eq}}}
\newcommand{\sel}{\mathsf{C}}
\newcommand{\gate}{\mathsf{G}}
\newcommand{\FV}{\operatorname{FV}}
\newcommand{\depth}{\operatorname{depth}}
\newcommand{\asm}{\operatorname{Asm}}
\newcommand{\self}{\operatorname{Self}}
\newcommand{\one}{S0}
\newcommand{\entB}{\vdash_{\Bprf}}
\newcommand{\entC}{\vdash_{\Cprf}}
\newcommand{\eqvE}{\mathrel{\simeq_E}}
\newcommand{\EqE}{\mathrel{\equiv_E}}
\newcommand{\repo}{\url{https://github.com/VladimirReshetnikov/ProveIt}}
\lstset{basicstyle=\ttfamily\footnotesize,breaklines=true,columns=fullflexible,
 keepspaces=true,showstringspaces=false,frame=single,rulecolor=\color{black!20}}
\title{\textbf{A one-body context calculus for\newline intuitionistic primitive recursive arithmetic}\\[6pt]
 \large Exact equivalence, complete rules, and measured implementation costs}
\author{Research prepared for Vladimir Reshetnikov}
\date{November 11, 2026}

\begin{document}
\maketitle
\begin{abstract}
The raw clause presentation of ConciseIPRA retains pairs
$[A;y.B(y)]$, with ordinary reading $A\to\forall y B(y)$.
We remove the independent antecedent field: every retained entry is just
a body $B(y)$, read $B(0)\to\forall yB(y)$. Ordinary assumptions are
represented by a pointwise iterator gate. The resulting calculus uses the
same raw terms, the same equation/$D$ formulas, atomic reflexivity, the
existing two-case evaluation rule, and three logical rules. We prove that
it has exactly the same open consequences as the established native
calculus, under an explicit representation of assumptions. Consequently
it is bi-interpretable and equiconsistent with full first-order iPRA, with
whole numerical domains, actual equality, and both numerical composites
given by $x=y$.

The size result is deliberately qualified. The complete new checker is
3,130 normalized bytes, versus 2,999 for a matched two-field control.
Thus the context grammar is simpler, but this implementation does not
improve the repository's proof-checker-size criterion. A separate,
semantics-preserving implementation refinement of the existing full Eq/$D$
checker reduces 788 bytes to 778, while also reducing its token and AST
counts; a 777-byte variant also included adds two AST nodes to the baseline.
These refinements introduce no mathematical change.
This package therefore does \emph{not} claim a new globally smallest
axiomatization or a mathematical size breakthrough.
\end{abstract}

\section{Result and comparison boundary}

The comparison snapshot is ProveIt commit
\begin{center}\small\texttt{11814771b7490dfbea0ed5fcf0f5b23586334764}.\end{center}
The inspected directory is
\texttt{Logic/PrimitiveRecursiveArithmetic/ConciseIPRA}.
Its README describes seventy-three research packages, multiple concise
calculi, and the Lean bridge from the native scalar language to iPRA
\cite{readme,nativebridge}. Several familiar reductions are already present:
one-term transport, equality from iteration, clause contexts, zero and full
normalization, and the fusion of equality with witness discharge. Those
are not claimed as discoveries of this report.

The new observation concerns the \emph{origin of contexts}. In a raw
clause proof starting from ordinary assumptions, a retained clause arises
in only two ways:
\[
 \asm(A)=[0=0;v.A],\qquad
 \self(B)=[B(0);y.B(y)].
\]
The first represents an assumption; the second comes from opening a
paired formula. We show that the first can also be represented by an
entry of the second kind. The primitive context datatype can therefore
lose its independent antecedent field entirely.

\begin{theorem}[Main equivalence, stated informally]
For Eq/$D$ formulas $A_1,\ldots,A_n,F$, the new calculus derives $F$ from
the retained bodies $\gate A_1,\ldots,\gate A_n$ if and only if the
established native calculus derives $F$ from ordinary assumptions
$A_1,\ldots,A_n$.
\end{theorem}

The theorem is about arbitrary finite proofs and open consequences. It is
stronger than agreement on equations in the standard natural numbers.
Its proof does not require a normalizer, a predecessor construction, a
new cut theorem, or unrestricted substitution into witness-selected
logical derivations.

\paragraph{What remains unmet.}
Under the earlier working criterion---the size of a complete standalone
proof checker---the new calculus has not beaten the current proposals.
The positive mathematical result is the deletion of one context field and
the corresponding exact equivalence theorem. The separate 10-byte
implementation improvement concerns the existing calculus. Combining
these into one purported new record would be incorrect.

\section{Complete presentation on raw syntax}
\label{sec:rules}

\subsection{Terms and formulas}

There is one numerical sort. Terms are finite expressions
\[
 t ::= x\mid 0\mid S(t)\mid I(n,a;u.h).
\]
Only $u$ in the last child $h$ is bound by $I$. Terms are \emph{raw}:
$I(0,a;u.h)$ and $a$ are different syntax trees. All instances below use
capture-avoiding substitution; syntactic equality is up to consistent
renaming of bound variables, or literal equality in the delivered
de Bruijn representation.

The complete formula grammar is
\[
 A ::= (s=t)\mid D_{x,y}A.
\]
Both $x$ and $y$ are bound in the matrix of $D$. There are no primitive
propositional connectives, quantifiers other than $D$, or separate $Q$
formulas. A formula is a finite $D$ prefix ending in a single equation.
The designated contradiction is $0=S0$.

The \emph{ordinary reading} of the native binder is
\[
 \bigl(D_{x,y}B\bigr)^\circ
 =\exists x\bigl(B^\circ(x,0)\to\forall y\,B^\circ(x,y)\bigr).
 \tag{D}
\]
This is an interpretation into ordinary intuitionistic first-order logic,
not an additional primitive inference rule.

\subsection{One-body contexts and the gate}

A context entry is a single formula body $B(v)$ with one distinguished
bound numerical coordinate $v$. Its meaning is fixed:
\[
 \mathcal Q_v B:=B^\circ(0)\to\forall v\,B^\circ(v).
 \tag{Context}
\]
A context is a finite list of such bodies. The binder belonging to an
entry is distinct from all $D$ binders inside that entry. There is no
independent antecedent field and no primitive ordinary-assumption entry.

Put
\[
 \sel(v,a,b):=I(v,a;u.b),\qquad u\notin\FV(b).
\]
For $A=D_{x_1,y_1}\cdots D_{x_d,y_d}(s=t)$ and a fresh coordinate $v$,
define the abbreviation
\[
 \gate_v A :=
 D_{x_1,y_1}\cdots D_{x_d,y_d}
 \bigl(\sel(v,0,s)=\sel(v,0,t)\bigr).
 \tag{Gate}
\]
The fresh iterator variables are unused. Every old bound coordinate keeps
its binding. In particular, $v$ is placed outside all $2d$ logical
coordinates, while the endpoint terms are lifted once when placed under
the new iterator binder. Gate is a construction on the existing grammar,
not a new formula constructor. For a formula family $A(v)$, the same
pointwise construction defines $\gate_v A(v)$; the coordinate must stay
fresh for its internal binders. In its use to represent an ordinary
assumption, $v$ is additionally absent from $A$.

\subsection{All primitive rules}

Write $\Gamma\entB A$ for the least derivability relation generated by
the following five rule families. Eval has two cases, so there are six
displayed inference schemata in total.

\paragraph{R: atomic reflexivity.}
For any well-scoped numerical term $t$,
\[
                    \frac{}{\Gamma\entB t=t}.\tag{R}
\]
Reflexivity under a $D$ prefix is derived later; it is not part of R.

\paragraph{Eval: computational equality.}
\[
 \frac{\Gamma\entB a=t\qquad
       \Gamma\entB I(0,a;u.h)=r}
      {\Gamma\entB r=t},
 \tag{$E_0$}
\]
\[
 \frac{\Gamma\entB I(c,a;u.h)=t\qquad
       \Gamma\entB I(Sc,a;u.h)=r}
      {\Gamma\entB r=h[t/u]}.
 \tag{$E_S$}
\]
There is no independent root-joining, congruence, symmetry, transitivity,
or computation axiom. Both premises of Eval are literal equations.

\paragraph{Use.}
If $B(v)$ is a retained context occurrence and $t$ is a term,
\[
        \frac{\Gamma\entB B(0)}{\Gamma\entB B(t)}.
        \tag{Use}
\]

\paragraph{Open.}
\[
 \frac{\Gamma\entB D_{x,y}B(x,y)\qquad
       \Gamma,B(b,y)\entB A}
      {\Gamma\entB A}.
 \tag{Open}
\]
The second premise adds the single body $B(b,y)$. The eigenparameter $b$
is fresh for $\Gamma$, $A$, and all passive parameters of $B$.

\paragraph{Intro.}
\[
 \frac{\Gamma,\gate_v\bigl(B(w,\eta)\bigr)
                    \entB B(w,Sz)}
      {\Gamma\entB D_{x,y}B(x,y)},
 \tag{Intro}
\]
where $z$ is fresh for $\Gamma$, $w$, and the passive parameters of $B$,
and
\[
 \eta=
 \begin{cases}
 z,&\text{if $B$ is a literal equation and $w$ is literally $0$},\\
 0,&\text{otherwise}.
 \end{cases}
 \tag{Mode}
\]
The distinguished gate coordinate $v$ is fresh and separate from $z$.
The mode test is on syntax, before any equation reasoning. It is the
existing witness-selected introduction, with its temporary assumption
represented by a gate body.

These rules and the finite syntax are the entire presentation. In
particular, ordinary intuitionistic logic, numerical case distinction,
equation induction, weakening, ordinary hypothesis access, and cut are
not additional primitive rules.

\section{Equality is recovered before the logical argument}
\label{sec:equality}

The following calculations use only R, $E_0$, and $E_S$. They work in
every retained context and are the arithmetic interface already used by
the repository's two-case Eval calculus \cite{eval}.

\begin{lemma}[Computation and joining]
Both iterator computation equations are derivable. Moreover, from
$a=b$ and $a=c$ one can derive $c=b$.
\end{lemma}
\begin{proof}
For zero computation, apply $E_0$ to $a=a$ and
$I(0,a;u.h)=I(0,a;u.h)$. For successor computation, apply $E_S$ to the
reflexivities of the preceding and current iterates.

For joining put $W=I(0,a;u.u)$. Apply $E_0$ to $a=c$ and $W=W$ to obtain
$W=c$. Apply $E_0$ again to $a=b$ and $W=c$ to obtain $c=b$.
\end{proof}

\begin{lemma}[Crossed one-term transport]
From $a=b$ and $h(a)=r$ one can derive $r=h(b)$.
\end{lemma}
\begin{proof}
Joining and R first give symmetry and transitivity. Put
\[
 W=I(0,a;u.h(u)),\qquad J=I(S0,a;u.h(u)).
\]
Zero Eval gives $W=a$. Successor Eval applied directly to $W=a$ and
$J=J$ gives $J=h(a)$. This step does not use congruence. Joining and
transitivity give $W=b$ and $J=r$ from the two supplied equalities.
Successor Eval applied to these gives $r=h(b)$.
\end{proof}

Consequently the pure subsystem provides ordinary equational reasoning
in every term template, including repeated occurrences and nested
iterator binders. Write $s\EqE t$ if this subsystem proves $s=t$, with
the displayed numerical parameters free. These pure derivations are
stable under capture-avoiding numerical substitution: their rules have
no witness-dependent mode test.

In particular,
\[
 \sel(0,0,s)\EqE 0,\qquad
 \sel(Sz,0,s)\EqE s.
 \tag{Endpoints}
\]
The second identity holds because the update ignores its iterator
coordinate. It does not require evaluating the preceding iterate.

\section{Gated assumptions are accessible without circularity}
\label{sec:access}

\begin{lemma}[Derived deep reflexivity]
Every well-scoped $D$ prefix above an equation $t=t$ is derivable in the
new calculus, in any context.
\end{lemma}
\begin{proof}
Induct externally on the number of $D$ nodes. At depth zero use R. At a
node $D_{x,y}B$, apply Intro with witness $1=S0$. This selects case mode.
The successor instance of $B$ again ends in a reflexive equation, so the
induction hypothesis proves it in the extended context, ignoring the
new gate body. No access to that body is needed.
\end{proof}

Weakening and injective renaming of free parameters are admissible by
direct induction on derivations, after freshening eigenparameters.
Injective renaming preserves the distinction between literal zero and a
variable. We do not assume that arbitrary substitution preserves an
Intro proof tree: it can change a nonzero witness into literal zero.

For formulas with the same alpha-aligned $D$ prefix, define
$A\eqvE A'$ when their corresponding left endpoints, and their
corresponding right endpoints, are related by $\EqE$. Bound coordinates
in these endpoint equations are temporarily treated as free numerical
parameters. This definition does \emph{not} presuppose formula transport.

For every depth $d$ we establish two assertions in this order:
\begin{description}
\item[$L_d$.] If $A\eqvE A'$ and both have depth $d$, a derivation of
either formula in any retained-body context can be converted into a
derivation of the other.
\item[$H_d$.] If $A$ has depth $d$ and $\gate_v A$ is a retained body,
then $\Gamma\entB A$.
\end{description}

\begin{lemma}[Conversion and access]
All assertions $L_d$ and $H_d$ hold.
\end{lemma}
\begin{proof}
At depth zero $L_0$ is equational congruence from
Section~\ref{sec:equality}.

Suppose $L_{d-1}$ and $H_{d-1}$ have been established. Let
$A=D_{x,y}B$ and $A'=D_{x,y}B'$. Open a supplied proof of $A$ at fresh
$b$, retaining the body $B(b,y)$. Introduce $A'$ with witness $b$.
A fresh variable is not literally zero, so this is case mode, even when
the matrix is an equation. The inner context now contains
$\gate_v B'(b,0)$.

Use $H_{d-1}$ to obtain $B'(b,0)$. Apply $L_{d-1}$, instantiated at
$x=b,y=0$, to obtain $B(b,0)$. Use the retained body $B(b,y)$ at $Sz$ to
obtain $B(b,Sz)$. Apply $L_{d-1}$ to obtain $B'(b,Sz)$. Intro proves
$A'$, and Open discharges $b$. This proves $L_d$ in one direction; the
other uses the symmetric endpoint equations.

Having established $L_d$, let $T_d$ denote the same $D$ prefix above
$0=0$. Derived deep reflexivity proves $T_d$. The endpoint equations give
\[
 \gate_0 A\eqvE T_d,\qquad \gate_1 A\eqvE A.
\]
Apply $L_d$ to get $\gate_0 A$. Use the retained body $\gate_v A$ at
$1$, obtaining $\gate_1 A$, and use $L_d$ once more to obtain $A$.
This proves $H_d$.
\end{proof}

The dependency order is explicit:
\[
                  L_0,\ H_0,\ L_1,\ H_1,\ L_2,\ H_2,\ldots.
\]
Conversion at depth $d$ uses assumption access only at depth $d-1$.
The proof is an external induction on finite formula syntax. It neither
adds quantified induction to arithmetic nor presupposes the hypothesis
access that it is proving.

\paragraph{Why bound coordinates matter.}
One cannot replace this argument by a casual appeal to one global
equality substitution. The leaf terms may depend on coordinates bound
inside the $D$ prefix, so the endpoint equations live underneath those
binders. The depth argument lifts those pointwise equations correctly.

\section{Exact translation from the established clause calculus}
\label{sec:complete}

The source calculus $\Cprf$ is the raw Eq/$D$ clause calculus formalized
as \texttt{CPrf} in \texttt{NativeClausal.lean} \cite{clausal}. It has
R, the same two Eval cases, Use, Open, and the same witness-selected
Intro. Its context entries are pairs $[A;y.B]$, read
$A\to\forall yB$. Source Use requires a derivation of the stored
antecedent and returns a chosen instance of the stored body.

Consider a source derivation beginning with ordinary assumptions
$\asm(A_i)=[0=0;v.A_i]$. Mark every context \emph{occurrence} by origin:
\[
 \begin{array}{c|c|c}
 \text{origin}&\text{source clause}&\text{target body}\\\hline
 \text{ordinary assumption}&[0=0;v.A]&\gate_v A\\
 \text{opening a paired formula}&[B(b,0);y.B(b,y)]&B(b,y).
 \end{array}
\]
Different occurrences can have different origins even when their source
clauses happen to be identical. These annotations belong to the proof
translation only; the new checker stores neither origin tags nor pairs.

\begin{theorem}[Constructive completeness translation]
For Eq/$D$ formulas $A_1,\ldots,A_n,F$,
\[
 \asm(A_1),\ldots,\asm(A_n)\entC F
 \quad\Longrightarrow\quad
 \gate A_1,\ldots,\gate A_n\entB F.
 \tag{C-to-B}
\]
\end{theorem}
\begin{proof}
Induct on the source proof, carrying the occurrence annotations.
R and Eval translate unchanged. If Use selects an ordinary-assumption
occurrence, replace that step by $H_{\depth(A)}$; its supplied proof of
$0=0$ need not be used. If Use selects an Open occurrence, its antecedent
is exactly $B(b,0)$, so the translated step is Use.

Source Open appends precisely an entry of the second kind, which the
target Open appends as its single body. Source Intro appends precisely
an ordinary assumption $B(w,\eta)$, whose translation is exactly the
gate body added by target Intro. Fresh numerical parameters and the
mode test are unchanged. Translate both branch derivations recursively.
\end{proof}

This argument needs no target cut theorem and no unrestricted
substitution into target proof trees. It also does not assert that every
arbitrary two-field clause is literally one self-based clause. The
invariant concerns the origins reachable in the proofs needed for all
ordinary sequents, which is sufficient for the theorem.

\section{Soundness and the bi-interpretation}
\label{sec:bi}

\subsection{The ordinary scalar intermediary}

The established intermediary $\Teq$ is intuitionistic first-order logic
with equality over the same scalar terms. Its nonlogical basis consists
of the two iterator computation equations, numerical coverage
\[
 \forall x\,(x=0\ \lor\ \exists y\,x=Sy),
\]
separation $0=S0\to\bot$, and induction for literal equations with
side assumptions. The pinned project proves its bi-interpretation with
iPRA, including both internal round trips \cite{scalarbridge}.
We use that result as an established upstream theorem, not as a new
formalization supplied here.

\begin{lemma}[Meaning of a gate context]
In $\Teq$, for every Eq/$D$ formula $A$,
\[
                  \mathcal Q_v(\gate_v A)\ \longleftrightarrow\ A^\circ.
                  \tag{Assumption}
\]
\end{lemma}
\begin{proof}
The zero endpoint is a true reflexive $D$ prefix, and every successor
endpoint is equivalent to $A^\circ$, by computation and the ordinary
logical congruence of the reading of $D$.
From the left side, supply its true zero antecedent and specialize at
$1$ to obtain $A^\circ$. Conversely, from $A^\circ$, numerical coverage
proves the gate at every argument: at zero use its true zero endpoint;
at a successor use $A^\circ$. This gives the displayed closure.
\end{proof}

\begin{theorem}[Soundness in the theory, not just the standard model]
If $\Gamma\entB A$, then $\Teq$ derives $A^\circ$ from the ordinary
readings $\mathcal Q B$ of the bodies in $\Gamma$.
\end{theorem}
\begin{proof}
Induct on the derivation. R is equality reflexivity; Eval follows from
computation and equality. Use is implication elimination followed by
universal specialization. Open is existential elimination for the
ordinary reading of $D$, with the stated eigenparameter condition.

For Intro, the added body has reading equivalent to $B(w,\eta)^\circ$
by the preceding lemma. In case mode the translated premise proves
$B(w,Sz)$ from $B(w,0)$; coverage yields all instances, and existential
introduction supplies witness $w$. In progress mode the matrix is a
literal equation and the premise is exactly its progressive step.
Literal-equation induction gives the conditional universal closure,
again with witness $w=0$. All passive assumptions remain in place.
\end{proof}

Combining soundness, the gate equivalence, the upstream native/$\Teq$
completeness and native formula round-trip theorems, the source proof
translation, and \texttt{cprf\_iff\_nprf\_of\_qfree} gives the precise
main result. The native/$\Teq$ bridge is needed in the direction from
ordinary soundness back to native derivability.

\begin{theorem}[Exact ordinary-sequent equivalence]
For Eq/$D$ formulas $A_1,\ldots,A_n,F$,
\[
 \gate A_1,\ldots,\gate A_n\entB F
 \quad\Longleftrightarrow\quad
 A_1,\ldots,A_n\vdash_{\Nprf}F.
 \tag{Equivalence}
\]
Here the right side is the pinned native scalar calculus, on its
$Q$-free fragment.
\end{theorem}

\subsection{Both directions on terms and formulas}

The equivalence just proved is the identity on raw terms and object
Eq/$D$ formulas. Only the representation of assumptions changes, with
its internal equivalence already proved. Eliminate a native $Q_yB$ by
$D_{x,y}B$ with $x$ fresh and unused; the numerical domain is inhabited
by $0$. This is the project's \texttt{unQ} translation \cite{clausal}.

For clarity, the inherited numerical maps have the following form.
In the scalar-to-iPRA direction, each fixed scalar iterator is compiled
to a primitive recursive definition with its seed and passive variables
as parameters. In the reverse direction, zero, successor and projections
are immediate, and composition is substitution. A primitive recursion
\[
 f(0,\vec a)=g(\vec a),\quad
 f(Sn,\vec a)=h(n,f(n,\vec a),\vec a)
\]
is compiled with the already-derived scalar pairing and projections as
\begin{align*}
 f(n,\vec a)=\pi_2 I\bigl(n,\langle0,g(\vec a)\rangle;u.
 \langle S\pi_1u, h(\pi_1u,\pi_2u,\vec a)\rangle\bigr).
\end{align*}
Here $g$ and $h$ on the scalar right side denote their already-compiled
scalar translations. The pairing, predecessor, reflection of equality, and correctness of
this compiler are inherited from the upstream scalar bridge. We do not
assume them in proving the new context theorem.

Let $\tau$ be the established scalar-to-PR term map and $\sigma$ its
inverse interpretation. Define the interpretation from the new calculus
by the ordinary reading followed by $\tau$. Define the reverse map by
the upstream logical-macro translation after $\sigma$, then \texttt{unQ}.
The upstream theorems \texttt{native\_roundTrip} and
\texttt{pra\_native\_roundTrip}, transported through (Equivalence), prove
that every formula is internally equivalent to its round trip.

Both numerical domains are whole; equality is actual equality. The
composite numerical interpretations have graph $x=y$, as in the
upstream \texttt{BiInterp} theorem. No restriction to standard numbers,
quotient domain, external truth test, or merely equiconsistent surrogate
is introduced by this construction.

\begin{corollary}
The new calculus is bi-interpretable with full first-order iPRA and
\[
         \varnothing\entB 0=S0
         \quad\Longleftrightarrow\quad
         \ipra\vdash\bot.
\]
\end{corollary}

\section{The entire ordinary logical interface}

This section makes the interpretation of arbitrary ordinary formulas
explicit. The macros are derived operations, not new rules of the
checker. Their validity follows from the established native logical
interface and (Equivalence) \cite{logic}.

First, $\mathcal Q_yB$ can be written in the object language as
$D_{x,y}B$ with unused $x$. Define a depth-raising envelope
\[
                  L(A):=D_{x,y}\gate_y A,\qquad x\notin\FV(A).
\]
Its ordinary reading is equivalent to $A$ by (Assumption). Repeated
$L$-padding can therefore align the $D$ depths of any two formulas.
After aligning and alpha-renaming their prefixes, define
\[
 \operatorname{Sel}_v(A,B)=
 D^{d}\bigl(\sel(v,s_A,s_B)=\sel(v,t_A,t_B)\bigr).
\]
It has endpoint $A$ at zero and endpoint $B$ at every successor.
The notation $D^d$ here abbreviates the common, explicitly bound prefix.

\begin{center}
\begin{tabular}{>{\(}l<{\)} >{\(}p{0.71\linewidth}<{\)}}
\toprule
\text{Ordinary connective}&\text{Native macro}\\
\midrule
\bot&0=S0\\[3pt]
A\to B&\mathcal Q_y\operatorname{Sel}_y(A,B)\\[3pt]
A\land B&\mathcal Q_y\gate_y\operatorname{Sel}_{\operatorname{par}(y)}(A,B)\\[3pt]
\exists x\,A(x)&D_{x,y}\gate_y A(x)\\[3pt]
A\lor B&\exists v\,\operatorname{Sel}_v(A,B)\\[3pt]
\forall x\,A(x)&A(0)\land\mathcal Q_y A(y)\\
\bottomrule
\end{tabular}
\end{center}
Freshness is understood in every macro. The displayed existential,
conjunction and universal symbols on right sides are expanded using the
earlier rows. The parity term can be taken as
\[
 \operatorname{par}(n)=I(n,0;u.\sel(u,1,0)).
\]
Its needed arithmetic properties are derived in the native theory.
Implication works because its zero selector instance is $A$, and its
successor instances are $B$. The conjunction gate has a true zero base;
its values at $1$ and $2$ supply the two conjuncts. Existential
quantification works because the gated base is true, making the
conditional closure equivalent to the matrix at the selected witness.

\paragraph{A padding trap.}
Plain vacuous $D$ padding would be incorrect. If both coordinates are
unused, $D_{x,y}A$ reads $\exists x(A\to A)$ and is always true.
The envelope $L(A)$ above is essential. This distinction matters when
combining formulas of different depths.

\section{Checker implementation and evidence}
\label{sec:implementation}

\subsection{Public contract}

The new reference checker accepts a JSON object with exactly the keys
\texttt{goal} and \texttt{proof}. The goal must be a closed well-formed
Eq/$D$ formula. Successful verification prints \texttt{ACCEPT}; parsing,
typing, scoping, shape or proof mismatches reject. Like the matched
control, its practical execution is subject to host memory and recursion
limits; the mathematical checker is the corresponding finite recursion
without a fixed resource cap.

The codec uses $0$ for zero, negative integers for fresh free-parameter
levels, positive integers for bound indices, \texttt{[0,t]} for successor,
and \texttt{[n,a,h]} for iteration. Equation and paired formulas are
\texttt{[s,t]} and \texttt{[B]}. In a $D$ matrix, $y$ has index $1$
and $x$ index $2$. Exact integer tests exclude JSON booleans.

The five proof packet tags are \texttt{r}, \texttt{v}, \texttt{u},
\texttt{e}, and \texttt{i}. Their detailed data fields and runnable
examples are described in the accompanying README. All contexts are
constructed by the verifier; the input cannot inject an arbitrary
retained body. Both equality premises of Eval must evaluate to equations.
Fresh levels are chosen from the scope counter and cannot escape an
Open or Intro branch into its outer conclusion.

For a formula of depth $d$, the gate's external clause coordinate is
placed at bound index $2d+1$. Each old endpoint is lifted once under
the new iterator update binder. This is the implementation point most
likely to be wrong in an informal ``wrap the endpoints'' translation.
The package includes an independent numerical-environment test of these
scope operations.

\subsection{Matched costs and separate refinements}

All principal sizes are computed from
\texttt{ast.unparse(ast.parse(source))}, followed by one LF byte, on
CPython 3.12.14.
Significant tokens exclude whitespace and structural tokenizer markers;
AST nodes count all nodes yielded by \texttt{ast.walk}. The executable
measurement script and exact sources are included.

\begin{center}
\begin{tabular}{p{0.50\linewidth}rrr}
\toprule
Complete checker & Tokens & AST & Bytes\\
\midrule
Matched raw two-field control & 848 & 986 & 2,999\\
New raw one-body calculus, atomic R & 901 & 1,041 & 3,130\\
Existing full Eq/$D$ compact checker & 358 & 413 & 788\\
Same existing calculus, preferred refinement & 357 & 412 & 778\\
Same existing calculus, byte-priority variant & 357 & 415 & 777\\
\bottomrule
\end{tabular}
\end{center}

The mathematical comparison adds about $4.37\%$ in normalized bytes.
The cost is the explicit gate construction and capture-avoiding lifting;
deleting the redundant context field does not pay for that machinery in
this implementation. A tested variant with primitive deep reflexivity
was larger still, at 3,208 bytes. Deep reflexivity is therefore kept as
a derived lemma in the main calculus.

The preferred 778-byte checker saves 10 bytes (about $1.27\%$), one
significant token and one AST node relative to the existing full Eq/$D$
checker. Its changes are equivalent validation expressions; its calculus
is unchanged. The byte-priority variant additionally rewrites a Boolean
dispatch and saves 11 bytes in total, but adds two AST nodes relative to
the baseline. It therefore does not improve all three coordinates.
The guarded compact 516-byte
codec reported by \texttt{canonical\_recursors} has a different public
validation contract. Its size must not be used as though it were the
full Eq/$D$ checker. The reported 256-token and 278-AST minima also
belong to other programs; combining coordinatewise minima into one
imaginary checker would be misleading.

\subsection{What was checked}

The package contains executable positive certificates, deliberately
invalid certificates, normalized source measurements, the independent
gate-scope test, and a source-level equivalence audit of the compact
implementation refinement. The final machine-readable evidence records
the exact checks and their counts. Reproduction commands are in the
README; scripts write their own evidence rather than relying on an
unverifiable narrative of a run. The executed suite passed 140 standalone
process checks: 34 raw-calculus certificates in both ordinary Python and
\texttt{python -O}, and 12 inherited native fixtures for each of three
checkers in both modes. It also passed 18,882 validator cases, producing
37,764 candidate/control comparisons for the two refined validators.

The positive certificates include genuine equation induction using its
gated hypothesis, ordinary Open/Use, both Eval cases, derived joining,
and nested-$D$ assumption extraction. In the last example the extracted
assumption is false in the standard model, so its branch cannot succeed
by ignoring that assumption. Invalid certificates exercise malformed
types, shape mismatches, indices, and capture attempts.

The focused independent gate test checks 36,462 cases, including $D$
depths zero through three, nested iterator binders, bound logical
coordinates, free parameters, and opening at numerical arguments zero,
one and two. Its evaluator uses explicit environments and does not
implement its expected result by calling the gate transformation again.

These tests are supplementary evidence. The new soundness,
completeness and bi-interpretation transfer are written metatheorems,
independently reviewed during this work. They have not been compiled
as new Lean or Rocq proofs. The inspected upstream theorem files are
included for provenance, but the upstream Lean project was not rebuilt
in this run. No standard-model test is presented as a proof of
bi-interpretability.

\section{Further research questions}

\begin{enumerate}
\item \textbf{Can one-body contexts reduce a complete checker?}
The present theorem removes an independent field, but explicit gates
cost more than that field. A useful next step must amortize gate lifting
or find an equally complete assumption representation with cheaper
scope operations. Any claim should include the whole validator.

\item \textbf{Can hypothesis access be built into a smaller uniform
inference without a hidden tag?}
The source-origin proof shows that only assumptions and self-based
closures are needed. Storing an extra origin tag is a simpler
implementation idea, but its measured direct variants did not improve
the control. A new inference discipline that avoids both tags and gates
would be more substantial.

\item \textbf{Can the new metatheorem be formalized with a small
dependency footprint?}
A Lean implementation should first formalize the pointwise gate and its
binding identities, then atomic equality, derived deep reflexivity, and
the simultaneous $L_d/H_d$ argument. Occurrence-origin annotations can
give completeness without first proving target cut.

\item \textbf{Can the introduction be made easier to substitute?}
The witness-selected mode is retained here. The repository already has
diagonal alternatives. Their combination with one-body contexts may
improve proof transformations, but would need a fresh matched cost
comparison and should not be counted as a new diagonal construction.

\item \textbf{What is a useful lower-bound question?}
There is no global minimality proof here. Restricted classes of
presentations---fixed term syntax, explicit primitive-recursive scope
operations, complete input validation and a stated language/runtime---may
support meaningful lower bounds or exhaustive comparisons. A byte count
alone, with an unspecified trusted interpreter, cannot establish an
absolute simplest axiomatization.
\end{enumerate}

\section*{Reproducibility and attribution}

All claims about the repository refer to the pinned commit, not to a
moving branch. The term language, $D$ reading, witness-selected mode,
Eval subsystem, source clause calculus and scalar/iPRA bridge are
inherited. The contribution proved here is their one-body context
presentation and the exact occurrence-origin translation.

The deliverable contains this source and compiled PDF, the complete new
checker, a matched control, the separately refined compact checker and
its baseline, verification scripts and certificates, source provenance,
and a checksum manifest. Reference source files remain unmodified and
are not advertised as a self-contained Lean project.

\begin{thebibliography}{9}
\raggedright
\bibitem{readme}
V. Reshetnikov, ProveIt, \emph{Concise presentations of intuitionistic
primitive recursive arithmetic}, README at the pinned commit.
\href{https://github.com/VladimirReshetnikov/ProveIt/blob/11814771b7490dfbea0ed5fcf0f5b23586334764/Logic/PrimitiveRecursiveArithmetic/ConciseIPRA/README.md}{Repository README}.

\bibitem{clausal}
ProveIt, \texttt{NativeClausal.lean}: \texttt{CPrf},
\texttt{cprf\_iff\_nprf\_of\_qfree} and
\texttt{cprf\_iff\_nprf}.
\href{https://github.com/VladimirReshetnikov/ProveIt/blob/11814771b7490dfbea0ed5fcf0f5b23586334764/Logic/PrimitiveRecursiveArithmetic/ConciseIPRA/Lean/ConciseIPRA/NativeClausal.lean}{Pinned source}.

\bibitem{eval}
ProveIt, \texttt{NativeEval.lean}: two-case computational equality and
\texttt{eprf\_iff\_nprf}.
\href{https://github.com/VladimirReshetnikov/ProveIt/blob/11814771b7490dfbea0ed5fcf0f5b23586334764/Logic/PrimitiveRecursiveArithmetic/ConciseIPRA/Lean/ConciseIPRA/NativeEval.lean}{Pinned source}.

\bibitem{logic}
ProveIt, \texttt{NativeLogic.lean}, \texttt{NativeComposite.lean}:
the ordinary logical interface and its native composites.
\href{https://github.com/VladimirReshetnikov/ProveIt/blob/11814771b7490dfbea0ed5fcf0f5b23586334764/Logic/PrimitiveRecursiveArithmetic/ConciseIPRA/Lean/ConciseIPRA/NativeLogic.lean}{Pinned logical rules}.

\bibitem{scalarbridge}
ProveIt, \texttt{BiInterpretation.lean}: the scalar/iPRA
\texttt{BiInterp}, theorem translations and round trips.
\href{https://github.com/VladimirReshetnikov/ProveIt/blob/11814771b7490dfbea0ed5fcf0f5b23586334764/Logic/PrimitiveRecursiveArithmetic/ConciseIPRA/Lean/ConciseIPRA/BiInterpretation.lean}{Pinned source}.

\bibitem{nativebridge}
ProveIt, \texttt{NativeBiInterpretation.lean}:
\texttt{native\_iff\_pra}, \texttt{pra\_iff\_native},
\texttt{native\_roundTrip}, \texttt{pra\_native\_roundTrip}, and
\texttt{native\_inconsistent\_iff}.
\href{https://github.com/VladimirReshetnikov/ProveIt/blob/11814771b7490dfbea0ed5fcf0f5b23586334764/Logic/PrimitiveRecursiveArithmetic/ConciseIPRA/Lean/ConciseIPRA/NativeBiInterpretation.lean}{Pinned source}.
\end{thebibliography}
\end{document}
